\pdfbookmark[1]{Abstract}{Abstract}
\chapter*{Abstract}

\textbf{The subject is the building block and the rule.} An atom is a point carrying values: a
count of protons, a count of valence electrons, a radius, an electronegativity. A molecule or a
solid is a set of such points in space. What tells that set from a random packing of the same atoms
is one rule, valence, and valence fixes a bond to a length before anyone measures it. The method
this research paper is written for reads any domain as points carrying values and asks one question of the
arrangement: how far it sits from the most disordered arrangement of its own parts. Chemistry is
that question with the atom as the part.

\textbf{What is quoted and what is new.} The building blocks and the rules in the first four
chapters are the settled content of general chemistry: the periodic law, the shell counts, the
octet, the electronegativity scale, the covalent and van der Waals radii, and the bond lengths
tabulated from crystallography and gas-phase spectroscopy. Nothing about them is claimed here. What
this research paper adds is their placement in the method, and the predictions that placement
makes. The atom is a representation, the octet is a screening condition, the disordered packing is a
reference, and the bond length is a test oracle of the strongest kind, a fact and not a consensus.

\textbf{Nothing in this research paper is measured here.} The method has representations for a
protein and for a crystal, and both already rest on chemistry: a backbone repeats at three because
three bond lengths are fixed, and a cell edge is a bond geometry tiled. It has no representation
for a molecule itself. This research paper states what such a reading would find, with the
predictions on the page before the reading exists.

\textbf{Why no set here is bounded.} A bound put on a quantity by judgment makes the reading report
the bound. Chemistry invites that error in a sharp form, because a bond either exists or it does
not and a tolerance is the obvious way to decide. This research paper keeps to the opposite
discipline: a departure is measured against a background drawn from the atoms themselves, and neither
threshold, cutoff radius nor bounded element set is chosen by judgment. The question asked is
always how far, never whether.
